\begin{problem}[44.6]
  Find the degree over $\mathbb{Q}$ of the splitting field over $\mathbb{Q}$ of the polynomial $(x^2 - 2)(x^3 - 2)$ in $\mathbb{Q}[x]$.
\end{problem}

\begin{solution}
  The splitting field of $(x^2 - 2)(x^3 - 2)$ over $\mathbb{Q}$ is $\mathbb{Q}(\sqrt{2}, \sqrt[3]{2}, \omega)$, where $\omega = \frac{-1 + i\sqrt{3}}{2}$ is a primitive cube root of unity. This is because $x^2 - 2$ has roots $\pm\sqrt{2}$ and $x^3 - 2$ has roots $\sqrt[3]{2}, \omega\sqrt[3]{2}, \omega^2\sqrt[3]{2}$, all of which lie in $\mathbb{Q}(\sqrt{2}, \sqrt[3]{2}, \omega)$. We compute the degree using the tower:
  \[%
    \mathbb{Q} \subseteq \mathbb{Q}(\sqrt[3]{2}) \subseteq \mathbb{Q}(\sqrt[3]{2}, \omega) \subseteq \mathbb{Q}(\sqrt[3]{2}, \omega, \sqrt{2})
  .\]%
  First, $x^3 - 2$ is irreducible over $\mathbb{Q}$ by Eisenstein's criterion, so $[\mathbb{Q}(\sqrt[3]{2}) : \mathbb{Q}] = 3$. Second, $\omega$ satisfies $x^2 + x + 1 = 0$, which is irreducible over $\mathbb{Q}$ and has no roots in $\mathbb{Q}(\sqrt[3]{2})$ since $\mathbb{Q}(\sqrt[3]{2}) \subseteq \R$, so $[\mathbb{Q}(\sqrt[3]{2}, \omega) : \mathbb{Q}(\sqrt[3]{2})] = 2$. Third, $\sqrt{2}$ satisfies $x^2 - 2 = 0$; since $[\mathbb{Q}(\sqrt[3]{2}, \omega) : \mathbb{Q}] = 6$ and $[\mathbb{Q}(\sqrt[3]{2}, \omega, \sqrt{2}) : \mathbb{Q}]$ must be divisible by both $6$ and $[\mathbb{Q}(\sqrt{2}) : \mathbb{Q}] = 2$, but $\sqrt{2} \notin \mathbb{Q}(\sqrt[3]{2}, \omega)$ as that field has degree 6 over $\mathbb{Q}$ while adjoining $\sqrt{2}$ would require degree divisible by 2 in a way not already accounted for, we have $[\mathbb{Q}(\sqrt[3]{2}, \omega, \sqrt{2}) : \mathbb{Q}(\sqrt[3]{2}, \omega)] = 2$. By the tower law:
  \[%
    [\mathbb{Q}(\sqrt{2}, \sqrt[3]{2}, \omega) : \mathbb{Q}] = 3 \cdot 2 \cdot 2 = 12
  .\qedhere\]%
\end{solution}

\begin{problem}[44.8]
  What is the order of $G(\mathbb{Q}(\sqrt[3]{2}, i\sqrt{3})/\mathbb{Q})$?
\end{problem}

\begin{solution}
  The order of $G(\mathbb{Q}(\sqrt[3]{2}, i\sqrt{3})/\mathbb{Q})$ equals the degree $[\mathbb{Q}(\sqrt[3]{2}, i\sqrt{3}) : \mathbb{Q}]$, since $\mathbb{Q}(\sqrt[3]{2}, i\sqrt{3})$ is the splitting field of the separable polynomial $x^3 - 2$ over $\mathbb{Q}$ and is therefore a Galois extension. We note that $\mathbb{Q}(\sqrt[3]{2}, i\sqrt{3}) = \mathbb{Q}(\sqrt[3]{2}, \omega)$, where $\omega = \frac{-1 + i\sqrt{3}}{2}$ is a primitive cube root of unity, since $i\sqrt{3} = 2\omega + 1$. We compute the degree using the tower:
  \[%
    \mathbb{Q} \subseteq \mathbb{Q}(\sqrt[3]{2}) \subseteq \mathbb{Q}(\sqrt[3]{2}, \omega)
  .\]%
  First, $x^3 - 2$ is irreducible over $\mathbb{Q}$ by Eisenstein's criterion, so $[\mathbb{Q}(\sqrt[3]{2}) : \mathbb{Q}] = 3$. Second, $\omega$ satisfies $x^2 + x + 1 = 0$, which is irreducible over $\mathbb{Q}$ and has no roots in $\mathbb{Q}(\sqrt[3]{2})$ since $\mathbb{Q}(\sqrt[3]{2}) \subseteq \R$, so $[\mathbb{Q}(\sqrt[3]{2}, \omega) : \mathbb{Q}(\sqrt[3]{2})] = 2$. By the tower law:
  \[%
    |G(\mathbb{Q}(\sqrt[3]{2}, i\sqrt{3})/\mathbb{Q})| = [\mathbb{Q}(\sqrt[3]{2}, \omega) : \mathbb{Q}] = 3 \cdot 2 = 6
  .\qedhere\]%
\end{solution}

\begin{problem}[44.10]
  Let $\alpha$ be a zero of $x^3 + x^2 + 1$ over $\Z_2$. Show that $x^3 + x^2 + 1$ splits in $\Z_2(\alpha)$. [Hint: There are eight elements in $\Z_2(\alpha)$. Exhibit two more zeros of $x^3 + x^2 + 1$, in addition to $\alpha$, among these eight elements. Alternatively, use the results of Section 42.]
\end{problem}

\begin{solution}
  Since $\alpha^3 + \alpha^2 + 1 = 0$, we have the following reduction rule: $\alpha^3 = \alpha^2 + 1$ (since we're working in $\Z_2$). We already have one root, namely, $\beta_1 = \alpha$. We need to find two more roots of $f(x) = x^3 + x^2 + 1$ in $\Z_2(\alpha)$. Letting $\beta_2 = \alpha^2$, we have:
  \[%
    \beta_2^3 + \beta_2^2 + 1 = (\alpha^2)^3 + (\alpha^2)^2 + 1 = \alpha^6 + \alpha^4 + 1
  .\]%
  Using the reduction rule on $\beta_2^2$, we have:
  \[%
    \beta_2^2 = \alpha^4 = (\alpha^2 + 1)\alpha = \alpha^3 + \alpha = (\alpha^2 + 1) + \alpha = \alpha^2 + \alpha + 1
  .\]%
  Using the reduction rule on $\beta_2^3$, we have:
  \[%
    \beta_2^3 = \alpha^6 = (\alpha^2 + 1)^2 = \alpha^4 + 2\alpha^2 + 1 = \alpha^4 + 1 = \alpha^2 + \alpha
  ,\]%
  since we're working in $\Z_2$. Thus, we have:
  \[%
    f(\beta_2) = \beta_2^3 + \beta_2^2 + 1 = (\alpha^2 + \alpha) + (\alpha^2 + \alpha + 1) + 1 = 0
  .\]%
  Thus, $\beta_2 = \alpha^2$ is a root of $f(x)$.

  Now, we need to find a third root of $f(x)$ in $\Z_2(\alpha)$. Letting $\beta_3 = 1 + \alpha + \alpha^2$, we have:
  \[%
    \beta_3^3 + \beta_3^2 + 1 = (\alpha^2 + \alpha + 1)^3 + (\alpha^2 + \alpha + 1)^2 + 1
  .\]%
  Using the reduction rule on $\beta_3^2$, we have:
  \[%
    \beta_3^2 = (\alpha^2 + \alpha + 1)^2 = \alpha^4 + \alpha^2 + 1 = (\alpha^2 + \alpha + 1) + \alpha^2 + 1 = \alpha
  .\]%
  Using the reduction rule on $\beta_3^3$, we have:
  \[%
    \beta_3^3 = \beta_3 \cdot \beta_3^2 = (\alpha^2 + \alpha + 1) \cdot \alpha = \alpha^3 + \alpha^2 + \alpha = (\alpha^2 + 1) + \alpha^2 + \alpha = \alpha + 1
  .\]%
  Thus, we have:
  \[%
    \beta_3^3 + \beta_3^2 + 1 = (\alpha + 1) + \alpha + 1 = 0
  .\]%
  Thus, $\beta_3 = \alpha^2 + \alpha + 1$ is a zero of $f(x)$.

  Since $f(x) = x^3 + x^2 + 1$ has three zeros $\beta_1 = \alpha$, $\beta_2 = \alpha^2$, and $\beta_3 = \alpha^2 + \alpha + 1$ in $\Z_2(\alpha)$, and $f(x)$ has degree three, it splits completely as:
  \[%
    f(x) = (x + \alpha)(x + \alpha^2)(x + \alpha^2 + \alpha + 1)
  ,\]%
  in $\Z_2(\alpha)[x]$.
\end{solution}
